%% ma: L10-B01-0029
%% lop: 10
%% chuong: 1
%% bai: 1
%% dang: 10.1.7
%% loai: TN4
%% muc: VD
%% dap_an: "C"
%% nguon: "Ảnh giáo viên gửi 10/10/2026 (ThS. Nguyễn Hoàng Việt, Chương 1 Mệnh đề), B-Đề 2, Phần 1 Câu 12"
%% kien_thuc: "Bài 1 (mệnh đề; suy luận logic)"
%% trang_thai: cho-duyet
%% ly_do: "Dạng toán mới đề xuất 10.1.7: suy luận logic từ các dữ kiện (xếp hàng, thang máy)"
%% ghi_chu: "Kiểm bằng Python: liệt kê đủ 12 cách xếp thỏa điều kiện, chỉ cặp (3,4) trong các phương án luôn khác giới. Lời giải gốc không có"
\begin{ex}
Một nhóm học sinh $M$, $N$, $P$, $Q$, $R$ xếp thành một hàng dọc trước một quầy nước giải khát. Dưới đây là các thông tin ghi nhận được từ các học sinh trên:
\begin{itemize}
\item $M$, $P$, $R$ là nam; $N$, $Q$ là nữ;
\item $N$ đứng ở vị trí thứ nhất hoặc thứ hai;
\item $M$ đứng trước $Q$;
\item Học sinh đứng sau cùng là nam.
\end{itemize}
Hai vị trí nào sau đây phải là hai học sinh khác giới tính?
\choice{Thứ hai và ba}{Thứ hai và năm}{\True Thứ ba và tư}{Thứ ba và năm}
\loigiai{Chỉ có $N,Q$ là nữ. Vị trí thứ ba và thứ tư không thể cùng là nữ vì $N$ đứng ở vị trí $1$ hoặc $2$. Chúng cũng không thể cùng là nam: khi đó $N,Q$ phải ở hai vị trí trong $\{1,2,5\}$, mà vị trí $5$ là nam nên $N,Q$ ở vị trí $1,2$; do $N$ ở vị trí $1$ hoặc $2$ nên $Q$ đứng trước $M$, trái giả thiết ``$M$ đứng trước $Q$''. Vậy hai vị trí thứ ba và thứ tư luôn khác giới tính. Các phương án khác không chắc: thứ hai và ba cùng giới với cách xếp $M,N,Q,P,R$; thứ hai và năm cùng giới với $N,M,P,Q,R$; thứ ba và năm cùng giới với $M,N,P,Q,R$.}
\end{ex}
